%%%%%%%%%%%%%%%%%%%%%%%%%%%%%%%%%%%%%%%%%%%%%%%%%%%%%%%%%%%%
% 9.3 DYNAMICAL SYSTEMS
% The Composition of Advanced Mathematics
%%%%%%%%%%%%%%%%%%%%%%%%%%%%%%%%%%%%%%%%%%%%%%%%%%%%%%%%%%%%

\section{Dynamical Systems}
\label{sec:dynamical-systems}

\prerequisite{
Ordinary differential equations, autonomous systems, linear algebra,
eigenvalues and eigenvectors, matrix exponentials, multivariable
calculus, Jacobian matrices, and basic numerical methods.
}

\learningobjectives{
By the end of this section, the reader should be able to:
\begin{itemize}
    \item define a dynamical system;
    \item distinguish continuous-time and discrete-time systems;
    \item define phase space, state, trajectory, and orbit;
    \item identify equilibrium points;
    \item analyze one-dimensional autonomous systems using phase lines;
    \item classify stable, unstable, and semistable equilibria;
    \item understand Lyapunov stability and asymptotic stability;
    \item distinguish local and global stability;
    \item analyze planar linear systems using eigenvalues;
    \item classify nodes, saddles, spirals, and centers;
    \item use trace and determinant to classify planar linear systems;
    \item linearize nonlinear systems near equilibria;
    \item understand hyperbolic equilibria;
    \item understand the Hartman--Grobman theorem conceptually;
    \item identify nullclines;
    \item sketch phase portraits;
    \item understand invariant sets;
    \item understand stable and unstable manifolds conceptually;
    \item understand Lyapunov functions;
    \item apply basic Lyapunov stability arguments;
    \item understand limit cycles;
    \item recognize conservative systems;
    \item analyze planar Hamiltonian-type systems;
    \item understand the Poincaré--Bendixson theorem conceptually;
    \item define bifurcation;
    \item analyze saddle-node bifurcations;
    \item analyze transcritical bifurcations;
    \item analyze pitchfork bifurcations;
    \item recognize Hopf bifurcations;
    \item understand bifurcation diagrams;
    \item analyze logistic-type nonlinear dynamics;
    \item understand predator--prey systems;
    \item understand competing-species systems conceptually;
    \item understand discrete dynamical systems;
    \item analyze fixed points of iterated maps;
    \item determine stability of fixed points using derivatives;
    \item understand period-two orbits;
    \item understand the logistic map;
    \item recognize period doubling and chaos;
    \item understand sensitivity to initial conditions;
    \item understand Lyapunov exponents conceptually;
    \item recognize strange attractors;
    \item distinguish deterministic chaos from randomness;
    \item understand numerical phase-plane exploration;
    \item connect dynamical systems with ODEs, difference equations,
          biology, physics, and nonlinear modeling.
\end{itemize}
}

%%%%%%%%%%%%%%%%%%%%%%%%%%%%%%%%%%%%%%%%%%%%%%%%%%%%%%%%%%%%
\subsection{What Is a Dynamical System?}
\label{subsec:what-is-dynamical-system}
%%%%%%%%%%%%%%%%%%%%%%%%%%%%%%%%%%%%%%%%%%%%%%%%%%%%%%%%%%%%

A dynamical system describes how a state changes over time according to a
specified rule.

For continuous time, a dynamical system often has the form

\[
\boxed{
\mathbf{x}'
=
\mathbf{f}(\mathbf{x}).
}
\]

For discrete time, it may have the form

\[
\boxed{
\mathbf{x}_{n+1}
=
\mathbf{F}(\mathbf{x}_n).
}
\]

The central questions are often qualitative:

\[
\boxed{
\text{Where do solutions go?}
}
\]

\[
\boxed{
\text{Which states are stable?}
}
\]

\[
\boxed{
\text{Do oscillations occur?}
}
\]

and

\[
\boxed{
\text{How does behavior change as parameters vary?}
}
\]

%%%%%%%%%%%%%%%%%%%%%%%%%%%%%%%%%%%%%%%%%%%%%%%%%%%%%%%%%%%%
\subsection{State of a System}
\label{subsec:dynamical-state}
%%%%%%%%%%%%%%%%%%%%%%%%%%%%%%%%%%%%%%%%%%%%%%%%%%%%%%%%%%%%

The state contains enough information to determine future evolution from
the model.

For a two-dimensional continuous system,

\[
\boxed{
\mathbf{x}(t)
=
\begin{pmatrix}
x(t)\\
y(t)
\end{pmatrix}.
}
\]

The ordered pair

\[
(x,y)
\]

represents the state at a particular time.

%%%%%%%%%%%%%%%%%%%%%%%%%%%%%%%%%%%%%%%%%%%%%%%%%%%%%%%%%%%%
\subsection{Phase Space}
\label{subsec:phase-space}
%%%%%%%%%%%%%%%%%%%%%%%%%%%%%%%%%%%%%%%%%%%%%%%%%%%%%%%%%%%%

\begin{definitionbox}{Phase Space}
The phase space is the set of all possible states of a dynamical system.
\end{definitionbox}

For a planar system,

\[
\boxed{
\text{phase space}
=
\mathbb{R}^2
}
\]

or some physically meaningful subset of it.

For an \(n\)-dimensional system,

\[
\boxed{
\mathbf{x}
\in
\mathbb{R}^n.
}
\]

%%%%%%%%%%%%%%%%%%%%%%%%%%%%%%%%%%%%%%%%%%%%%%%%%%%%%%%%%%%%
\subsection{Trajectory}
\label{subsec:trajectory-dynamical-system}
%%%%%%%%%%%%%%%%%%%%%%%%%%%%%%%%%%%%%%%%%%%%%%%%%%%%%%%%%%%%

A trajectory is the path traced by the state as time evolves.

For

\[
\mathbf{x}'
=
\mathbf{f}(\mathbf{x}),
\]

a solution

\[
\mathbf{x}(t)
\]

determines a trajectory in phase space.

The direction along the trajectory is determined by increasing time.

%%%%%%%%%%%%%%%%%%%%%%%%%%%%%%%%%%%%%%%%%%%%%%%%%%%%%%%%%%%%
\subsection{Orbit}
\label{subsec:orbit-dynamical-system}
%%%%%%%%%%%%%%%%%%%%%%%%%%%%%%%%%%%%%%%%%%%%%%%%%%%%%%%%%%%%

The orbit of a state is the set of states visited by its trajectory.

For continuous systems,

\[
\boxed{
\mathcal{O}(\mathbf{x}_0)
=
\{
\mathbf{x}(t):
t\in I
\}.
}
\]

The parameterization by time may be ignored when discussing only the
geometric orbit.

%%%%%%%%%%%%%%%%%%%%%%%%%%%%%%%%%%%%%%%%%%%%%%%%%%%%%%%%%%%%
\subsection{Flow}
\label{subsec:dynamical-flow}
%%%%%%%%%%%%%%%%%%%%%%%%%%%%%%%%%%%%%%%%%%%%%%%%%%%%%%%%%%%%

If the solution beginning at \(\mathbf{x}_0\) is denoted by

\[
\phi_t(\mathbf{x}_0),
\]

then the family

\[
\boxed{
\phi_t
}
\]

is called the flow.

For an autonomous system, the flow satisfies

\[
\boxed{
\phi_0(\mathbf{x})
=
\mathbf{x},
}
\]

and

\[
\boxed{
\phi_{t+s}(\mathbf{x})
=
\phi_t
(
\phi_s(\mathbf{x})
).
}
\]

%%%%%%%%%%%%%%%%%%%%%%%%%%%%%%%%%%%%%%%%%%%%%%%%%%%%%%%%%%%%
\subsection{Equilibrium Points}
\label{subsec:dynamical-equilibria}
%%%%%%%%%%%%%%%%%%%%%%%%%%%%%%%%%%%%%%%%%%%%%%%%%%%%%%%%%%%%

An equilibrium point

\[
\boxed{
\mathbf{x}^*
}
\]

satisfies

\[
\boxed{
\mathbf{f}(\mathbf{x}^*)
=
\mathbf{0}.
}
\]

If the system begins at \(\mathbf{x}^*\), it remains there for all time.

Equilibria are also called:

\[
\boxed{
\text{fixed points},
}
\]

\[
\boxed{
\text{critical points},
}
\]

or

\[
\boxed{
\text{steady states}.
}
\]

Terminology depends on context.

%%%%%%%%%%%%%%%%%%%%%%%%%%%%%%%%%%%%%%%%%%%%%%%%%%%%%%%%%%%%
\subsection{One-Dimensional Autonomous Systems}
\label{subsec:one-dimensional-dynamical-systems}
%%%%%%%%%%%%%%%%%%%%%%%%%%%%%%%%%%%%%%%%%%%%%%%%%%%%%%%%%%%%

Consider

\[
\boxed{
x'
=
f(x).
}
\]

Equilibria satisfy

\[
\boxed{
f(x^*)=0.
}
\]

Between equilibria:

\[
f(x)>0
\]

means \(x\) increases, while

\[
f(x)<0
\]

means \(x\) decreases.

%%%%%%%%%%%%%%%%%%%%%%%%%%%%%%%%%%%%%%%%%%%%%%%%%%%%%%%%%%%%
\subsection{Phase-Line Analysis}
\label{subsec:dynamical-phase-line}
%%%%%%%%%%%%%%%%%%%%%%%%%%%%%%%%%%%%%%%%%%%%%%%%%%%%%%%%%%%%

A phase line summarizes the direction of motion along the real axis.

For

\[
x'=f(x),
\]

mark all equilibria, then determine the sign of \(f(x)\) between them.

Use arrows:

\[
\boxed{
\rightarrow
\quad
\text{if }f(x)>0,
}
\]

and

\[
\boxed{
\leftarrow
\quad
\text{if }f(x)<0.
}
\]

%%%%%%%%%%%%%%%%%%%%%%%%%%%%%%%%%%%%%%%%%%%%%%%%%%%%%%%%%%%%
\subsection{Stable Equilibrium}
\label{subsec:dynamical-stable-equilibrium}
%%%%%%%%%%%%%%%%%%%%%%%%%%%%%%%%%%%%%%%%%%%%%%%%%%%%%%%%%%%%

An equilibrium is stable if sufficiently small perturbations remain
small.

Informally,

\[
\boxed{
\text{nearby initial states remain nearby}.
}
\]

For one-dimensional systems, arrows pointing toward an equilibrium from
both sides indicate asymptotic stability.

%%%%%%%%%%%%%%%%%%%%%%%%%%%%%%%%%%%%%%%%%%%%%%%%%%%%%%%%%%%%
\subsection{Unstable Equilibrium}
\label{subsec:dynamical-unstable-equilibrium}
%%%%%%%%%%%%%%%%%%%%%%%%%%%%%%%%%%%%%%%%%%%%%%%%%%%%%%%%%%%%

An equilibrium is unstable if nearby states can move away.

For a phase line, arrows pointing away on both sides indicate instability.

%%%%%%%%%%%%%%%%%%%%%%%%%%%%%%%%%%%%%%%%%%%%%%%%%%%%%%%%%%%%
\subsection{Semistable Equilibrium}
\label{subsec:semistable-equilibrium}
%%%%%%%%%%%%%%%%%%%%%%%%%%%%%%%%%%%%%%%%%%%%%%%%%%%%%%%%%%%%

An equilibrium can attract trajectories from one side and repel them from
the other.

Such an equilibrium is often called semistable.

It is not asymptotically stable in the ordinary two-sided sense.

%%%%%%%%%%%%%%%%%%%%%%%%%%%%%%%%%%%%%%%%%%%%%%%%%%%%%%%%%%%%
\subsection{Derivative Stability Criterion}
\label{subsec:dynamical-derivative-stability}
%%%%%%%%%%%%%%%%%%%%%%%%%%%%%%%%%%%%%%%%%%%%%%%%%%%%%%%%%%%%

Suppose

\[
x'=f(x)
\]

and

\[
f(x^*)=0.
\]

If

\[
\boxed{
f'(x^*)<0,
}
\]

then \(x^*\) is locally asymptotically stable.

If

\[
\boxed{
f'(x^*)>0,
}
\]

then \(x^*\) is unstable.

If

\[
\boxed{
f'(x^*)=0,
}
\]

the criterion is inconclusive.

%%%%%%%%%%%%%%%%%%%%%%%%%%%%%%%%%%%%%%%%%%%%%%%%%%%%%%%%%%%%
\subsection{Worked Example: Phase-Line Stability}
\label{subsec:worked-phase-line-stability}
%%%%%%%%%%%%%%%%%%%%%%%%%%%%%%%%%%%%%%%%%%%%%%%%%%%%%%%%%%%%

\begin{workedexamplebox}{One-Dimensional Dynamics}

Consider

\[
x'
=
x(1-x).
\]

Equilibria satisfy

\[
x(1-x)=0.
\]

Thus,

\[
x^*=0,
\qquad
x^*=1.
\]

For

\[
x<0,
\]

we have

\[
x'<0.
\]

For

\[
0<x<1,
\]

we have

\[
x'>0.
\]

For

\[
x>1,
\]

we have

\[
x'<0.
\]

Therefore,

\begin{answerbox}
\[
\boxed{
x=0
\text{ is unstable},
}
\]

and

\[
\boxed{
x=1
\text{ is asymptotically stable}.
}
\]
\end{answerbox}

\end{workedexamplebox}

%%%%%%%%%%%%%%%%%%%%%%%%%%%%%%%%%%%%%%%%%%%%%%%%%%%%%%%%%%%%
\subsection{Lyapunov Stability}
\label{subsec:lyapunov-stability-definition}
%%%%%%%%%%%%%%%%%%%%%%%%%%%%%%%%%%%%%%%%%%%%%%%%%%%%%%%%%%%%

\begin{definitionbox}{Lyapunov Stability}
An equilibrium \(\mathbf{x}^*\) is stable in the sense of Lyapunov if,
for every \(\varepsilon>0\), there exists \(\delta>0\) such that

\[
\|
\mathbf{x}(0)-\mathbf{x}^*
\|
<
\delta
\]

implies

\[
\boxed{
\|
\mathbf{x}(t)-\mathbf{x}^*
\|
<
\varepsilon
}
\]

for all \(t\geq0\).
\end{definitionbox}

This says small initial perturbations remain small.

%%%%%%%%%%%%%%%%%%%%%%%%%%%%%%%%%%%%%%%%%%%%%%%%%%%%%%%%%%%%
\subsection{Asymptotic Stability}
\label{subsec:asymptotic-stability}
%%%%%%%%%%%%%%%%%%%%%%%%%%%%%%%%%%%%%%%%%%%%%%%%%%%%%%%%%%%%

An equilibrium is asymptotically stable if it is Lyapunov stable and
nearby solutions satisfy

\[
\boxed{
\mathbf{x}(t)
\to
\mathbf{x}^*
\qquad
\text{as }
t\to\infty.
}
\]

Thus asymptotic stability combines:

\[
\boxed{
\text{staying nearby}
}
\]

and

\[
\boxed{
\text{eventual convergence}.
}
\]

%%%%%%%%%%%%%%%%%%%%%%%%%%%%%%%%%%%%%%%%%%%%%%%%%%%%%%%%%%%%
\subsection{Exponential Stability}
\label{subsec:exponential-stability}
%%%%%%%%%%%%%%%%%%%%%%%%%%%%%%%%%%%%%%%%%%%%%%%%%%%%%%%%%%%%

An equilibrium is exponentially stable if nearby perturbations satisfy a
bound of the form

\[
\boxed{
\|
\mathbf{x}(t)-\mathbf{x}^*
\|
\leq
M
e^{-\alpha t}
\|
\mathbf{x}(0)-\mathbf{x}^*
\|
}
\]

for some

\[
M\geq1,
\qquad
\alpha>0.
\]

This is stronger than ordinary asymptotic stability.

%%%%%%%%%%%%%%%%%%%%%%%%%%%%%%%%%%%%%%%%%%%%%%%%%%%%%%%%%%%%
\subsection{Local Versus Global Stability}
\label{subsec:local-global-stability}
%%%%%%%%%%%%%%%%%%%%%%%%%%%%%%%%%%%%%%%%%%%%%%%%%%%%%%%%%%%%

Local stability describes behavior only for initial conditions near an
equilibrium.

Global asymptotic stability means all admissible initial conditions
approach the equilibrium.

Thus:

\[
\boxed{
\text{local}
\rightarrow
\text{nearby behavior},
}
\]

\[
\boxed{
\text{global}
\rightarrow
\text{entire state space or domain}.
}
\]

%%%%%%%%%%%%%%%%%%%%%%%%%%%%%%%%%%%%%%%%%%%%%%%%%%%%%%%%%%%%
\subsection{Basin of Attraction}
\label{subsec:basin-attraction}
%%%%%%%%%%%%%%%%%%%%%%%%%%%%%%%%%%%%%%%%%%%%%%%%%%%%%%%%%%%%

The basin of attraction of an asymptotically stable equilibrium
\(\mathbf{x}^*\) is

\[
\boxed{
\mathcal{B}(\mathbf{x}^*)
=
\left\{
\mathbf{x}_0:
\phi_t(\mathbf{x}_0)
\to
\mathbf{x}^*
\right\}.
}
\]

Different attractors may have different basins.

%%%%%%%%%%%%%%%%%%%%%%%%%%%%%%%%%%%%%%%%%%%%%%%%%%%%%%%%%%%%
\subsection{Linear Planar Systems}
\label{subsec:dynamical-linear-planar-systems}
%%%%%%%%%%%%%%%%%%%%%%%%%%%%%%%%%%%%%%%%%%%%%%%%%%%%%%%%%%%%

Consider

\[
\boxed{
\mathbf{x}'
=
A\mathbf{x},
}
\]

where

\[
A
\in
\mathbb{R}^{2\times2}.
\]

The equilibrium

\[
\mathbf{x}=\mathbf{0}
\]

is classified largely by the eigenvalues of \(A\).

%%%%%%%%%%%%%%%%%%%%%%%%%%%%%%%%%%%%%%%%%%%%%%%%%%%%%%%%%%%%
\subsection{Stable Node}
\label{subsec:dynamical-stable-node}
%%%%%%%%%%%%%%%%%%%%%%%%%%%%%%%%%%%%%%%%%%%%%%%%%%%%%%%%%%%%

If the eigenvalues are real, negative, and independent,

\[
\boxed{
\lambda_1<0,
\qquad
\lambda_2<0,
}
\]

the origin is a stable node.

Solutions approach the origin exponentially.

%%%%%%%%%%%%%%%%%%%%%%%%%%%%%%%%%%%%%%%%%%%%%%%%%%%%%%%%%%%%
\subsection{Unstable Node}
\label{subsec:dynamical-unstable-node}
%%%%%%%%%%%%%%%%%%%%%%%%%%%%%%%%%%%%%%%%%%%%%%%%%%%%%%%%%%%%

If both eigenvalues are real and positive,

\[
\boxed{
\lambda_1>0,
\qquad
\lambda_2>0,
}
\]

the origin is an unstable node.

Trajectories move away from the equilibrium.

%%%%%%%%%%%%%%%%%%%%%%%%%%%%%%%%%%%%%%%%%%%%%%%%%%%%%%%%%%%%
\subsection{Saddle Point}
\label{subsec:dynamical-saddle}
%%%%%%%%%%%%%%%%%%%%%%%%%%%%%%%%%%%%%%%%%%%%%%%%%%%%%%%%%%%%

If the eigenvalues have opposite signs,

\[
\boxed{
\lambda_1\lambda_2<0,
}
\]

the equilibrium is a saddle.

One eigendirection is stable and the other is unstable.

Every saddle equilibrium is unstable.

%%%%%%%%%%%%%%%%%%%%%%%%%%%%%%%%%%%%%%%%%%%%%%%%%%%%%%%%%%%%
\subsection{Stable Spiral}
\label{subsec:stable-spiral}
%%%%%%%%%%%%%%%%%%%%%%%%%%%%%%%%%%%%%%%%%%%%%%%%%%%%%%%%%%%%

Suppose

\[
\lambda
=
\alpha
\pm
i\beta,
\]

with

\[
\beta\neq0.
\]

If

\[
\boxed{
\alpha<0,
}
\]

trajectories spiral toward the equilibrium.

The equilibrium is a stable spiral or spiral sink.

%%%%%%%%%%%%%%%%%%%%%%%%%%%%%%%%%%%%%%%%%%%%%%%%%%%%%%%%%%%%
\subsection{Unstable Spiral}
\label{subsec:unstable-spiral}
%%%%%%%%%%%%%%%%%%%%%%%%%%%%%%%%%%%%%%%%%%%%%%%%%%%%%%%%%%%%

If

\[
\boxed{
\alpha>0,
}
\]

trajectories spiral away.

The equilibrium is an unstable spiral or spiral source.

%%%%%%%%%%%%%%%%%%%%%%%%%%%%%%%%%%%%%%%%%%%%%%%%%%%%%%%%%%%%
\subsection{Center}
\label{subsec:dynamical-center}
%%%%%%%%%%%%%%%%%%%%%%%%%%%%%%%%%%%%%%%%%%%%%%%%%%%%%%%%%%%%

If the eigenvalues are purely imaginary,

\[
\boxed{
\lambda
=
\pm i\beta,
}
\]

a linear planar system can have closed orbits surrounding the equilibrium.

The origin is a center.

It is stable in the Lyapunov sense but not asymptotically stable.

%%%%%%%%%%%%%%%%%%%%%%%%%%%%%%%%%%%%%%%%%%%%%%%%%%%%%%%%%%%%
\subsection{Repeated Eigenvalues}
\label{subsec:dynamical-repeated-eigenvalues}
%%%%%%%%%%%%%%%%%%%%%%%%%%%%%%%%%%%%%%%%%%%%%%%%%%%%%%%%%%%%

Repeated eigenvalues require checking the number of independent
eigenvectors.

If the matrix has two independent eigenvectors, a star-type node may
occur.

If it is defective, generalized eigenvectors produce an improper node.

The sign of the repeated eigenvalue determines attraction or repulsion.

%%%%%%%%%%%%%%%%%%%%%%%%%%%%%%%%%%%%%%%%%%%%%%%%%%%%%%%%%%%%
\subsection{Trace and Determinant}
\label{subsec:dynamical-trace-determinant}
%%%%%%%%%%%%%%%%%%%%%%%%%%%%%%%%%%%%%%%%%%%%%%%%%%%%%%%%%%%%

For

\[
A
=
\begin{pmatrix}
a&b\\
c&d
\end{pmatrix},
\]

define

\[
\boxed{
\tau
=
\operatorname{tr}(A)
=
a+d,
}
\]

and

\[
\boxed{
\Delta
=
\det(A)
=
ad-bc.
}
\]

The eigenvalues satisfy

\[
\boxed{
\lambda^2
-
\tau\lambda
+
\Delta
=
0.
}
\]

%%%%%%%%%%%%%%%%%%%%%%%%%%%%%%%%%%%%%%%%%%%%%%%%%%%%%%%%%%%%
\subsection{Trace--Determinant Plane}
\label{subsec:trace-determinant-plane}
%%%%%%%%%%%%%%%%%%%%%%%%%%%%%%%%%%%%%%%%%%%%%%%%%%%%%%%%%%%%

The discriminant is

\[
\boxed{
D
=
\tau^2-4\Delta.
}
\]

Then:

\[
\boxed{
\Delta<0
\Rightarrow
\text{saddle},
}
\]

\[
\boxed{
\Delta>0,\ D>0
\Rightarrow
\text{node-type behavior},
}
\]

and

\[
\boxed{
\Delta>0,\ D<0
\Rightarrow
\text{spiral or center-type behavior}.
}
\]

The sign of \(\tau\) determines attraction or repulsion when real parts
are nonzero.

%%%%%%%%%%%%%%%%%%%%%%%%%%%%%%%%%%%%%%%%%%%%%%%%%%%%%%%%%%%%
\subsection{Worked Example: Linear Classification}
\label{subsec:worked-linear-classification}
%%%%%%%%%%%%%%%%%%%%%%%%%%%%%%%%%%%%%%%%%%%%%%%%%%%%%%%%%%%%

\begin{workedexamplebox}{Classifying an Equilibrium}

Consider

\[
\mathbf{x}'
=
\begin{pmatrix}
-2&0\\
0&-5
\end{pmatrix}
\mathbf{x}.
\]

The eigenvalues are

\[
\lambda_1=-2,
\qquad
\lambda_2=-5.
\]

Both are negative and real.

Therefore,

\begin{answerbox}
\[
\boxed{
\mathbf{0}
\text{ is a stable node}.
}
\]
\end{answerbox}

\end{workedexamplebox}

%%%%%%%%%%%%%%%%%%%%%%%%%%%%%%%%%%%%%%%%%%%%%%%%%%%%%%%%%%%%
\subsection{Nonlinear Planar Systems}
\label{subsec:nonlinear-planar-systems}
%%%%%%%%%%%%%%%%%%%%%%%%%%%%%%%%%%%%%%%%%%%%%%%%%%%%%%%%%%%%

A nonlinear planar autonomous system has form

\[
\boxed{
\begin{cases}
x'=f(x,y),\\
y'=g(x,y).
\end{cases}
}
\]

Equilibria satisfy

\[
\boxed{
f(x,y)=0,
\qquad
g(x,y)=0.
}
\]

Exact solutions are rarely necessary for qualitative analysis.

%%%%%%%%%%%%%%%%%%%%%%%%%%%%%%%%%%%%%%%%%%%%%%%%%%%%%%%%%%%%
\subsection{Nullclines}
\label{subsec:nullclines}
%%%%%%%%%%%%%%%%%%%%%%%%%%%%%%%%%%%%%%%%%%%%%%%%%%%%%%%%%%%%

The \(x\)-nullcline is the set where

\[
\boxed{
x'=f(x,y)=0.
}
\]

The \(y\)-nullcline is the set where

\[
\boxed{
y'=g(x,y)=0.
}
\]

Their intersections are equilibria.

Nullclines divide the phase plane into regions where signs of \(x'\) and
\(y'\) are constant or predictable.

%%%%%%%%%%%%%%%%%%%%%%%%%%%%%%%%%%%%%%%%%%%%%%%%%%%%%%%%%%%%
\subsection{Direction from Nullclines}
\label{subsec:nullcline-direction}
%%%%%%%%%%%%%%%%%%%%%%%%%%%%%%%%%%%%%%%%%%%%%%%%%%%%%%%%%%%%

In each region:

\[
x'>0
\]

means motion to the right,

\[
x'<0
\]

means motion to the left,

\[
y'>0
\]

means motion upward,

and

\[
y'<0
\]

means motion downward.

Combining these signs gives the local flow direction.

%%%%%%%%%%%%%%%%%%%%%%%%%%%%%%%%%%%%%%%%%%%%%%%%%%%%%%%%%%%%
\subsection{Jacobian Matrix}
\label{subsec:dynamical-jacobian}
%%%%%%%%%%%%%%%%%%%%%%%%%%%%%%%%%%%%%%%%%%%%%%%%%%%%%%%%%%%%

For

\[
\mathbf{f}(x,y)
=
\begin{pmatrix}
f(x,y)\\
g(x,y)
\end{pmatrix},
\]

the Jacobian is

\[
\boxed{
J(x,y)
=
\begin{pmatrix}
f_x&f_y\\
g_x&g_y
\end{pmatrix}.
}
\]

At equilibrium

\[
(x^*,y^*),
\]

the Jacobian gives the linearization.

%%%%%%%%%%%%%%%%%%%%%%%%%%%%%%%%%%%%%%%%%%%%%%%%%%%%%%%%%%%%
\subsection{Linearization Near Equilibrium}
\label{subsec:dynamical-linearization}
%%%%%%%%%%%%%%%%%%%%%%%%%%%%%%%%%%%%%%%%%%%%%%%%%%%%%%%%%%%%

Let

\[
\mathbf{x}
=
\mathbf{x}^*
+
\mathbf{u}.
\]

Then near an equilibrium,

\[
\mathbf{u}'
=
\mathbf{f}
(
\mathbf{x}^*+\mathbf{u}
)
\]

is approximated by

\[
\boxed{
\mathbf{u}'
\approx
J(\mathbf{x}^*)
\mathbf{u}.
}
\]

Thus eigenvalues of the Jacobian often determine local behavior.

%%%%%%%%%%%%%%%%%%%%%%%%%%%%%%%%%%%%%%%%%%%%%%%%%%%%%%%%%%%%
\subsection{Hyperbolic Equilibria}
\label{subsec:dynamical-hyperbolic-equilibria}
%%%%%%%%%%%%%%%%%%%%%%%%%%%%%%%%%%%%%%%%%%%%%%%%%%%%%%%%%%%%

\begin{definitionbox}{Hyperbolic Equilibrium}
An equilibrium is hyperbolic if no eigenvalue of its Jacobian has zero
real part.
\end{definitionbox}

For hyperbolic equilibria, linearization reliably determines the local
qualitative stability type.

%%%%%%%%%%%%%%%%%%%%%%%%%%%%%%%%%%%%%%%%%%%%%%%%%%%%%%%%%%%%
\subsection{Hartman--Grobman Theorem Preview}
\label{subsec:hartman-grobman-preview}
%%%%%%%%%%%%%%%%%%%%%%%%%%%%%%%%%%%%%%%%%%%%%%%%%%%%%%%%%%%%

Near a hyperbolic equilibrium, a nonlinear system is locally
topologically equivalent to its linearization.

Thus the nonlinear and linearized phase portraits have the same local
qualitative structure.

This is the content of the Hartman--Grobman theorem.

%%%%%%%%%%%%%%%%%%%%%%%%%%%%%%%%%%%%%%%%%%%%%%%%%%%%%%%%%%%%
\subsection{When Linearization Is Inconclusive}
\label{subsec:linearization-inconclusive}
%%%%%%%%%%%%%%%%%%%%%%%%%%%%%%%%%%%%%%%%%%%%%%%%%%%%%%%%%%%%

If the Jacobian has an eigenvalue with zero real part, linearization may
fail to determine stability.

Examples include:

\[
\boxed{
\lambda=0,
}
\]

or

\[
\boxed{
\lambda=\pm i\omega.
}
\]

Nonlinear terms must then be analyzed more carefully.

%%%%%%%%%%%%%%%%%%%%%%%%%%%%%%%%%%%%%%%%%%%%%%%%%%%%%%%%%%%%
\subsection{Invariant Sets}
\label{subsec:invariant-sets}
%%%%%%%%%%%%%%%%%%%%%%%%%%%%%%%%%%%%%%%%%%%%%%%%%%%%%%%%%%%%

A set \(S\) is invariant if trajectories beginning in \(S\) remain in
\(S\).

Symbolically,

\[
\boxed{
\mathbf{x}_0\in S
\Rightarrow
\phi_t(\mathbf{x}_0)\in S
}
\]

for all times for which the solution is defined.

Equilibria and periodic orbits are examples of invariant sets.

%%%%%%%%%%%%%%%%%%%%%%%%%%%%%%%%%%%%%%%%%%%%%%%%%%%%%%%%%%%%
\subsection{Stable Manifold Preview}
\label{subsec:stable-manifold-preview}
%%%%%%%%%%%%%%%%%%%%%%%%%%%%%%%%%%%%%%%%%%%%%%%%%%%%%%%%%%%%

For a hyperbolic equilibrium, the stable manifold consists of initial
states that approach the equilibrium as

\[
t\to\infty.
\]

Near the equilibrium, it is tangent to the eigenspace associated with
eigenvalues having negative real part.

%%%%%%%%%%%%%%%%%%%%%%%%%%%%%%%%%%%%%%%%%%%%%%%%%%%%%%%%%%%%
\subsection{Unstable Manifold Preview}
\label{subsec:unstable-manifold-preview}
%%%%%%%%%%%%%%%%%%%%%%%%%%%%%%%%%%%%%%%%%%%%%%%%%%%%%%%%%%%%

The unstable manifold consists of states that approach the equilibrium as
time runs backward.

Near the equilibrium, it is tangent to eigendirections corresponding to
eigenvalues with positive real part.

For a saddle, stable and unstable manifolds organize the phase portrait.

%%%%%%%%%%%%%%%%%%%%%%%%%%%%%%%%%%%%%%%%%%%%%%%%%%%%%%%%%%%%
\subsection{Lyapunov Functions}
\label{subsec:lyapunov-functions}
%%%%%%%%%%%%%%%%%%%%%%%%%%%%%%%%%%%%%%%%%%%%%%%%%%%%%%%%%%%%

A Lyapunov function is an energy-like scalar function used to study
stability without solving the differential equation explicitly.

Let

\[
V(\mathbf{x})
\]

satisfy

\[
\boxed{
V(\mathbf{x}^*)=0
}
\]

and

\[
\boxed{
V(\mathbf{x})>0
}
\]

near \(\mathbf{x}^*\) for

\[
\mathbf{x}\neq\mathbf{x}^*.
\]

%%%%%%%%%%%%%%%%%%%%%%%%%%%%%%%%%%%%%%%%%%%%%%%%%%%%%%%%%%%%
\subsection{Derivative Along Trajectories}
\label{subsec:lyapunov-derivative}
%%%%%%%%%%%%%%%%%%%%%%%%%%%%%%%%%%%%%%%%%%%%%%%%%%%%%%%%%%%%

For

\[
\mathbf{x}'
=
\mathbf{f}(\mathbf{x}),
\]

the derivative of \(V\) along trajectories is

\[
\boxed{
\dot V
=
\nabla V(\mathbf{x})^T
\mathbf{f}(\mathbf{x}).
}
\]

If

\[
\dot V\leq0,
\]

then \(V\) does not increase along solutions.

%%%%%%%%%%%%%%%%%%%%%%%%%%%%%%%%%%%%%%%%%%%%%%%%%%%%%%%%%%%%
\subsection{Basic Lyapunov Stability Criterion}
\label{subsec:basic-lyapunov-criterion}
%%%%%%%%%%%%%%%%%%%%%%%%%%%%%%%%%%%%%%%%%%%%%%%%%%%%%%%%%%%%

If \(V\) is positive definite near an equilibrium and

\[
\boxed{
\dot V\leq0,
}
\]

then the equilibrium can often be shown to be stable.

If

\[
\boxed{
\dot V<0
}
\]

away from the equilibrium, asymptotic stability can often be concluded.

Precise statements require appropriate regularity and domain conditions.

%%%%%%%%%%%%%%%%%%%%%%%%%%%%%%%%%%%%%%%%%%%%%%%%%%%%%%%%%%%%
\subsection{Worked Example: Lyapunov Function}
\label{subsec:worked-lyapunov}
%%%%%%%%%%%%%%%%%%%%%%%%%%%%%%%%%%%%%%%%%%%%%%%%%%%%%%%%%%%%

\begin{workedexamplebox}{Quadratic Lyapunov Function}

Consider

\[
x'=-x,
\qquad
y'=-2y.
\]

Choose

\[
V(x,y)
=
x^2+y^2.
\]

Then

\[
\dot V
=
2xx'
+
2yy'.
\]

Thus,

\[
\dot V
=
-2x^2
-
4y^2.
\]

Therefore,

\begin{answerbox}
\[
\boxed{
\dot V<0
\quad
\text{for }
(x,y)\neq(0,0).
}
\]

Hence the origin is asymptotically stable.
\end{answerbox}

\end{workedexamplebox}

%%%%%%%%%%%%%%%%%%%%%%%%%%%%%%%%%%%%%%%%%%%%%%%%%%%%%%%%%%%%
\subsection{LaSalle Invariance Principle Preview}
\label{subsec:lasalle-preview}
%%%%%%%%%%%%%%%%%%%%%%%%%%%%%%%%%%%%%%%%%%%%%%%%%%%%%%%%%%%%

When

\[
\dot V\leq0
\]

but is not strictly negative, Lyapunov's basic criterion may not establish
asymptotic stability.

LaSalle's invariance principle studies the largest invariant set
contained in

\[
\boxed{
\{
\mathbf{x}:
\dot V(\mathbf{x})=0
\}.
}
\]

This often proves convergence even when \(\dot V\) is only negative
semidefinite.

%%%%%%%%%%%%%%%%%%%%%%%%%%%%%%%%%%%%%%%%%%%%%%%%%%%%%%%%%%%%
\subsection{Periodic Orbits}
\label{subsec:periodic-orbits}
%%%%%%%%%%%%%%%%%%%%%%%%%%%%%%%%%%%%%%%%%%%%%%%%%%%%%%%%%%%%

A nonconstant solution is periodic if there exists \(T>0\) such that

\[
\boxed{
\mathbf{x}(t+T)
=
\mathbf{x}(t)
}
\]

for all \(t\).

The smallest positive such \(T\), when it exists, is the period.

Its trajectory forms a closed orbit in phase space.

%%%%%%%%%%%%%%%%%%%%%%%%%%%%%%%%%%%%%%%%%%%%%%%%%%%%%%%%%%%%
\subsection{Limit Cycles}
\label{subsec:limit-cycles}
%%%%%%%%%%%%%%%%%%%%%%%%%%%%%%%%%%%%%%%%%%%%%%%%%%%%%%%%%%%%

A limit cycle is an isolated periodic orbit.

Nearby trajectories may:

\[
\boxed{
\text{approach it},
}
\]

\[
\boxed{
\text{move away from it},
}
\]

or exhibit different behavior on opposite sides.

A stable limit cycle represents sustained oscillation.

%%%%%%%%%%%%%%%%%%%%%%%%%%%%%%%%%%%%%%%%%%%%%%%%%%%%%%%%%%%%
\subsection{Stable Limit Cycle}
\label{subsec:stable-limit-cycle}
%%%%%%%%%%%%%%%%%%%%%%%%%%%%%%%%%%%%%%%%%%%%%%%%%%%%%%%%%%%%

A periodic orbit is asymptotically stable if nearby trajectories approach
the orbit as

\[
t\to\infty.
\]

The state continues moving around the cycle rather than converging to one
point.

Thus the attractor is a closed orbit.

%%%%%%%%%%%%%%%%%%%%%%%%%%%%%%%%%%%%%%%%%%%%%%%%%%%%%%%%%%%%
\subsection{Conservative Systems}
\label{subsec:conservative-systems}
%%%%%%%%%%%%%%%%%%%%%%%%%%%%%%%%%%%%%%%%%%%%%%%%%%%%%%%%%%%%

A conservative mechanical system often has a conserved energy

\[
\boxed{
E(x,y)
=
\text{constant along trajectories}.
}
\]

Then

\[
\boxed{
\frac{dE}{dt}=0.
}
\]

Trajectories lie on level sets

\[
\boxed{
E(x,y)=C.
}
\]

%%%%%%%%%%%%%%%%%%%%%%%%%%%%%%%%%%%%%%%%%%%%%%%%%%%%%%%%%%%%
\subsection{Second-Order Conservative Equation}
\label{subsec:second-order-conservative}
%%%%%%%%%%%%%%%%%%%%%%%%%%%%%%%%%%%%%%%%%%%%%%%%%%%%%%%%%%%%

Consider

\[
\boxed{
x''
+
V'(x)
=
0.
}
\]

Set

\[
y=x'.
\]

Then

\[
\boxed{
\begin{cases}
x'=y,\\
y'=-V'(x).
\end{cases}
}
\]

The energy is

\[
\boxed{
E(x,y)
=
\frac12y^2+V(x).
}
\]

%%%%%%%%%%%%%%%%%%%%%%%%%%%%%%%%%%%%%%%%%%%%%%%%%%%%%%%%%%%%
\subsection{Energy Conservation}
\label{subsec:energy-conservation-dynamics}
%%%%%%%%%%%%%%%%%%%%%%%%%%%%%%%%%%%%%%%%%%%%%%%%%%%%%%%%%%%%

Differentiate:

\[
\frac{dE}{dt}
=
yy'
+
V'(x)x'.
\]

Using

\[
x'=y,
\qquad
y'=-V'(x),
\]

we obtain

\[
\boxed{
\frac{dE}{dt}
=
-yV'(x)+V'(x)y
=
0.
}
\]

Therefore energy is constant along trajectories.

%%%%%%%%%%%%%%%%%%%%%%%%%%%%%%%%%%%%%%%%%%%%%%%%%%%%%%%%%%%%
\subsection{Simple Pendulum}
\label{subsec:simple-pendulum-dynamics}
%%%%%%%%%%%%%%%%%%%%%%%%%%%%%%%%%%%%%%%%%%%%%%%%%%%%%%%%%%%%

The undamped pendulum satisfies

\[
\boxed{
\theta''
+
\frac{g}{L}
\sin\theta
=
0.
}
\]

Set

\[
\omega=\theta'.
\]

Then

\[
\boxed{
\begin{cases}
\theta'=\omega,\\
\omega'=
-\dfrac{g}{L}\sin\theta.
\end{cases}
}
\]

This is nonlinear because of \(\sin\theta\).

%%%%%%%%%%%%%%%%%%%%%%%%%%%%%%%%%%%%%%%%%%%%%%%%%%%%%%%%%%%%
\subsection{Pendulum Energy}
\label{subsec:pendulum-energy}
%%%%%%%%%%%%%%%%%%%%%%%%%%%%%%%%%%%%%%%%%%%%%%%%%%%%%%%%%%%%

A conserved energy is

\[
\boxed{
E
=
\frac12
\omega^2
+
\frac{g}{L}
(
1-\cos\theta
).
}
\]

Different energy levels correspond to:

\[
\boxed{
\text{oscillation},
}
\]

\[
\boxed{
\text{separatrix motion},
}
\]

or

\[
\boxed{
\text{continuous rotation}.
}
\]

%%%%%%%%%%%%%%%%%%%%%%%%%%%%%%%%%%%%%%%%%%%%%%%%%%%%%%%%%%%%
\subsection{Poincaré--Bendixson Theorem Preview}
\label{subsec:poincare-bendixson-preview}
%%%%%%%%%%%%%%%%%%%%%%%%%%%%%%%%%%%%%%%%%%%%%%%%%%%%%%%%%%%%

Planar continuous autonomous systems have important restrictions on their
long-term behavior.

Roughly, if a bounded trajectory remains in a region containing only
finitely many equilibria and its omega-limit set contains no equilibrium,
then under suitable hypotheses the limiting set is a periodic orbit.

This is a central idea of the Poincaré--Bendixson theorem.

%%%%%%%%%%%%%%%%%%%%%%%%%%%%%%%%%%%%%%%%%%%%%%%%%%%%%%%%%%%%
\subsection{No Continuous-Time Chaos in Planar Autonomous ODEs}
\label{subsec:no-planar-continuous-chaos}
%%%%%%%%%%%%%%%%%%%%%%%%%%%%%%%%%%%%%%%%%%%%%%%%%%%%%%%%%%%%

The Poincaré--Bendixson theory implies that smooth autonomous ODE systems
in two dimensions cannot exhibit the same type of deterministic chaotic
attractors found in higher-dimensional continuous systems.

Continuous-time chaos generally requires at least three state variables.

%%%%%%%%%%%%%%%%%%%%%%%%%%%%%%%%%%%%%%%%%%%%%%%%%%%%%%%%%%%%
\subsection{Predator--Prey Model}
\label{subsec:predator-prey}
%%%%%%%%%%%%%%%%%%%%%%%%%%%%%%%%%%%%%%%%%%%%%%%%%%%%%%%%%%%%

A classical predator--prey model is

\[
\boxed{
\begin{cases}
x'
=
ax-bxy,
\\
y'
=
-cy+dxy,
\end{cases}
}
\]

where

\[
x
=
\text{prey population},
\]

and

\[
y
=
\text{predator population}.
\]

The constants

\[
a,b,c,d
\]

are positive.

%%%%%%%%%%%%%%%%%%%%%%%%%%%%%%%%%%%%%%%%%%%%%%%%%%%%%%%%%%%%
\subsection{Predator--Prey Equilibria}
\label{subsec:predator-prey-equilibria}
%%%%%%%%%%%%%%%%%%%%%%%%%%%%%%%%%%%%%%%%%%%%%%%%%%%%%%%%%%%%

Set

\[
x'=0,
\qquad
y'=0.
\]

Since

\[
x'
=
x(a-by),
\]

and

\[
y'
=
y(-c+dx),
\]

the equilibria are

\[
\boxed{
(0,0)
}
\]

and

\[
\boxed{
\left(
\frac{c}{d},
\frac{a}{b}
\right).
}
\]

%%%%%%%%%%%%%%%%%%%%%%%%%%%%%%%%%%%%%%%%%%%%%%%%%%%%%%%%%%%%
\subsection{Predator--Prey Nullclines}
\label{subsec:predator-prey-nullclines}
%%%%%%%%%%%%%%%%%%%%%%%%%%%%%%%%%%%%%%%%%%%%%%%%%%%%%%%%%%%%

The prey nullclines are

\[
\boxed{
x=0
}
\]

and

\[
\boxed{
y=\frac{a}{b}.
}
\]

The predator nullclines are

\[
\boxed{
y=0
}
\]

and

\[
\boxed{
x=\frac{c}{d}.
}
\]

These lines divide the positive quadrant into regions of increasing and
decreasing populations.

%%%%%%%%%%%%%%%%%%%%%%%%%%%%%%%%%%%%%%%%%%%%%%%%%%%%%%%%%%%%
\subsection{Competition Model Preview}
\label{subsec:competition-model-preview}
%%%%%%%%%%%%%%%%%%%%%%%%%%%%%%%%%%%%%%%%%%%%%%%%%%%%%%%%%%%%

A two-species competition model may have form

\[
\boxed{
x'
=
r_1x
\left(
1-\frac{x+\alpha y}{K_1}
\right),
}
\]

\[
\boxed{
y'
=
r_2y
\left(
1-\frac{y+\beta x}{K_2}
\right).
}
\]

Depending on parameter values, the model can exhibit:

\[
\boxed{
\text{competitive exclusion},
}
\]

\[
\boxed{
\text{coexistence},
}
\]

or

\[
\boxed{
\text{bistability}.
}
\]

%%%%%%%%%%%%%%%%%%%%%%%%%%%%%%%%%%%%%%%%%%%%%%%%%%%%%%%%%%%%
\subsection{Bifurcations}
\label{subsec:bifurcations}
%%%%%%%%%%%%%%%%%%%%%%%%%%%%%%%%%%%%%%%%%%%%%%%%%%%%%%%%%%%%

A bifurcation occurs when the qualitative behavior of a dynamical system
changes as a parameter varies.

Consider

\[
\boxed{
x'
=
f(x,\mu),
}
\]

where

\[
\mu
\]

is a parameter.

Changes may include:

\[
\boxed{
\text{creation of equilibria},
}
\]

\[
\boxed{
\text{loss of stability},
}
\]

or

\[
\boxed{
\text{birth of periodic orbits}.
}
\]

%%%%%%%%%%%%%%%%%%%%%%%%%%%%%%%%%%%%%%%%%%%%%%%%%%%%%%%%%%%%
\subsection{Bifurcation Diagram}
\label{subsec:bifurcation-diagram}
%%%%%%%%%%%%%%%%%%%%%%%%%%%%%%%%%%%%%%%%%%%%%%%%%%%%%%%%%%%%

A bifurcation diagram plots equilibrium values against a parameter.

Typically:

\[
\boxed{
\text{solid branch}
=
\text{stable equilibrium},
}
\]

and

\[
\boxed{
\text{dashed branch}
=
\text{unstable equilibrium}.
}
\]

The convention should be stated whenever used.

%%%%%%%%%%%%%%%%%%%%%%%%%%%%%%%%%%%%%%%%%%%%%%%%%%%%%%%%%%%%
\subsection{Saddle-Node Bifurcation}
\label{subsec:saddle-node-bifurcation}
%%%%%%%%%%%%%%%%%%%%%%%%%%%%%%%%%%%%%%%%%%%%%%%%%%%%%%%%%%%%

The normal form is

\[
\boxed{
x'
=
\mu-x^2.
}
\]

Equilibria satisfy

\[
x^2=\mu.
\]

Therefore:

\[
\mu<0
\Rightarrow
\text{no equilibria},
\]

\[
\mu=0
\Rightarrow
x=0,
\]

and

\[
\mu>0
\Rightarrow
x=\pm\sqrt{\mu}.
\]

A pair of equilibria is created or destroyed at the bifurcation.

%%%%%%%%%%%%%%%%%%%%%%%%%%%%%%%%%%%%%%%%%%%%%%%%%%%%%%%%%%%%
\subsection{Stability in the Saddle-Node Normal Form}
\label{subsec:saddle-node-stability}
%%%%%%%%%%%%%%%%%%%%%%%%%%%%%%%%%%%%%%%%%%%%%%%%%%%%%%%%%%%%

For

\[
f(x)=\mu-x^2,
\]

\[
f'(x)=-2x.
\]

At

\[
x=+\sqrt{\mu},
\]

\[
f'<0,
\]

so the equilibrium is stable.

At

\[
x=-\sqrt{\mu},
\]

\[
f'>0,
\]

so it is unstable.

%%%%%%%%%%%%%%%%%%%%%%%%%%%%%%%%%%%%%%%%%%%%%%%%%%%%%%%%%%%%
\subsection{Transcritical Bifurcation}
\label{subsec:transcritical-bifurcation}
%%%%%%%%%%%%%%%%%%%%%%%%%%%%%%%%%%%%%%%%%%%%%%%%%%%%%%%%%%%%

A standard normal form is

\[
\boxed{
x'
=
\mu x-x^2.
}
\]

Factor:

\[
\boxed{
x'
=
x(\mu-x).
}
\]

The equilibria are

\[
\boxed{
x=0,
\qquad
x=\mu.
}
\]

At

\[
\mu=0,
\]

the branches intersect and exchange stability.

%%%%%%%%%%%%%%%%%%%%%%%%%%%%%%%%%%%%%%%%%%%%%%%%%%%%%%%%%%%%
\subsection{Pitchfork Bifurcation}
\label{subsec:pitchfork-bifurcation}
%%%%%%%%%%%%%%%%%%%%%%%%%%%%%%%%%%%%%%%%%%%%%%%%%%%%%%%%%%%%

A supercritical pitchfork normal form is

\[
\boxed{
x'
=
\mu x-x^3.
}
\]

Equilibria satisfy

\[
x
(
\mu-x^2
)
=
0.
\]

Thus

\[
\boxed{
x=0
}
\]

for all \(\mu\), and when

\[
\mu>0,
\]

additional equilibria appear:

\[
\boxed{
x=\pm\sqrt{\mu}.
}
\]

%%%%%%%%%%%%%%%%%%%%%%%%%%%%%%%%%%%%%%%%%%%%%%%%%%%%%%%%%%%%
\subsection{Supercritical Pitchfork Stability}
\label{subsec:supercritical-pitchfork-stability}
%%%%%%%%%%%%%%%%%%%%%%%%%%%%%%%%%%%%%%%%%%%%%%%%%%%%%%%%%%%%

Here,

\[
f'(x)
=
\mu-3x^2.
\]

At

\[
x=0,
\]

\[
f'(0)=\mu.
\]

Thus the origin is stable when

\[
\mu<0
\]

and unstable when

\[
\mu>0.
\]

At

\[
x=\pm\sqrt{\mu},
\]

\[
f'
=
\mu-3\mu
=
-2\mu<0.
\]

Thus the new branches are stable.

%%%%%%%%%%%%%%%%%%%%%%%%%%%%%%%%%%%%%%%%%%%%%%%%%%%%%%%%%%%%
\subsection{Hopf Bifurcation Preview}
\label{subsec:hopf-bifurcation-preview}
%%%%%%%%%%%%%%%%%%%%%%%%%%%%%%%%%%%%%%%%%%%%%%%%%%%%%%%%%%%%

A Hopf bifurcation occurs when a pair of complex-conjugate eigenvalues
crosses the imaginary axis as a parameter varies.

Schematically,

\[
\boxed{
\lambda(\mu)
=
\alpha(\mu)
\pm
i\omega(\mu),
}
\]

with

\[
\boxed{
\alpha(\mu_0)=0.
}
\]

Under suitable nondegeneracy conditions, a periodic orbit may be created
or destroyed.

%%%%%%%%%%%%%%%%%%%%%%%%%%%%%%%%%%%%%%%%%%%%%%%%%%%%%%%%%%%%
\subsection{Supercritical and Subcritical Hopf Preview}
\label{subsec:hopf-types-preview}
%%%%%%%%%%%%%%%%%%%%%%%%%%%%%%%%%%%%%%%%%%%%%%%%%%%%%%%%%%%%

A supercritical Hopf bifurcation often creates a small stable limit cycle.

A subcritical Hopf bifurcation often involves an unstable limit cycle and
can produce abrupt loss of stability.

The precise classification depends on nonlinear terms.

%%%%%%%%%%%%%%%%%%%%%%%%%%%%%%%%%%%%%%%%%%%%%%%%%%%%%%%%%%%%
\subsection{Discrete Dynamical Systems}
\label{subsec:discrete-dynamical-systems}
%%%%%%%%%%%%%%%%%%%%%%%%%%%%%%%%%%%%%%%%%%%%%%%%%%%%%%%%%%%%

A discrete dynamical system evolves by iteration:

\[
\boxed{
x_{n+1}
=
f(x_n).
}
\]

Starting from

\[
x_0,
\]

one obtains

\[
x_1=f(x_0),
\]

\[
x_2=f(x_1),
\]

and so on.

%%%%%%%%%%%%%%%%%%%%%%%%%%%%%%%%%%%%%%%%%%%%%%%%%%%%%%%%%%%%
\subsection{Iterates}
\label{subsec:dynamical-iterates}
%%%%%%%%%%%%%%%%%%%%%%%%%%%%%%%%%%%%%%%%%%%%%%%%%%%%%%%%%%%%

The second iterate is

\[
\boxed{
f^2(x)
=
f(f(x)).
}
\]

The \(n\)-th iterate is

\[
\boxed{
f^n(x)
=
\underbrace{
f\circ\cdots\circ f
}_{n\text{ times}}
(x).
}
\]

Then

\[
\boxed{
x_n
=
f^n(x_0).
}
\]

%%%%%%%%%%%%%%%%%%%%%%%%%%%%%%%%%%%%%%%%%%%%%%%%%%%%%%%%%%%%
\subsection{Fixed Points of a Map}
\label{subsec:fixed-points-map}
%%%%%%%%%%%%%%%%%%%%%%%%%%%%%%%%%%%%%%%%%%%%%%%%%%%%%%%%%%%%

A fixed point \(x^*\) satisfies

\[
\boxed{
f(x^*)=x^*.
}
\]

If

\[
x_0=x^*,
\]

then

\[
x_n=x^*
\]

for every \(n\).

%%%%%%%%%%%%%%%%%%%%%%%%%%%%%%%%%%%%%%%%%%%%%%%%%%%%%%%%%%%%
\subsection{Stability of a Fixed Point}
\label{subsec:fixed-point-stability-map}
%%%%%%%%%%%%%%%%%%%%%%%%%%%%%%%%%%%%%%%%%%%%%%%%%%%%%%%%%%%%

Suppose

\[
f(x^*)=x^*.
\]

If

\[
\boxed{
|f'(x^*)|<1,
}
\]

then the fixed point is locally attracting.

If

\[
\boxed{
|f'(x^*)|>1,
}
\]

it is locally repelling.

If

\[
|f'(x^*)|=1,
\]

the test is inconclusive.

%%%%%%%%%%%%%%%%%%%%%%%%%%%%%%%%%%%%%%%%%%%%%%%%%%%%%%%%%%%%
\subsection{Why the Derivative Criterion Works}
\label{subsec:fixed-point-derivative-reason}
%%%%%%%%%%%%%%%%%%%%%%%%%%%%%%%%%%%%%%%%%%%%%%%%%%%%%%%%%%%%

Let

\[
x_n=x^*+e_n,
\]

where \(e_n\) is small.

Using linear approximation,

\[
f(x^*+e_n)
\approx
f(x^*)
+
f'(x^*)e_n.
\]

Since

\[
f(x^*)=x^*,
\]

we obtain

\[
\boxed{
e_{n+1}
\approx
f'(x^*)e_n.
}
\]

Thus perturbations shrink when

\[
|f'(x^*)|<1.
\]

%%%%%%%%%%%%%%%%%%%%%%%%%%%%%%%%%%%%%%%%%%%%%%%%%%%%%%%%%%%%
\subsection{Cobweb Diagrams}
\label{subsec:cobweb-diagrams}
%%%%%%%%%%%%%%%%%%%%%%%%%%%%%%%%%%%%%%%%%%%%%%%%%%%%%%%%%%%%

For

\[
x_{n+1}=f(x_n),
\]

a cobweb diagram uses the curves

\[
\boxed{
y=f(x)
}
\]

and

\[
\boxed{
y=x.
}
\]

Alternating vertical and horizontal movements visualize iteration.

This gives a graphical method for studying convergence and oscillation.

%%%%%%%%%%%%%%%%%%%%%%%%%%%%%%%%%%%%%%%%%%%%%%%%%%%%%%%%%%%%
\subsection{Period-Two Orbits}
\label{subsec:period-two-orbits}
%%%%%%%%%%%%%%%%%%%%%%%%%%%%%%%%%%%%%%%%%%%%%%%%%%%%%%%%%%%%

A period-two orbit consists of distinct points

\[
x_1,
\qquad
x_2
\]

such that

\[
\boxed{
f(x_1)=x_2,
}
\]

and

\[
\boxed{
f(x_2)=x_1.
}
\]

Equivalently,

\[
\boxed{
f^2(x_1)=x_1,
}
\]

but

\[
\boxed{
f(x_1)\neq x_1.
}
\]

%%%%%%%%%%%%%%%%%%%%%%%%%%%%%%%%%%%%%%%%%%%%%%%%%%%%%%%%%%%%
\subsection{Stability of a Periodic Orbit}
\label{subsec:periodic-orbit-map-stability}
%%%%%%%%%%%%%%%%%%%%%%%%%%%%%%%%%%%%%%%%%%%%%%%%%%%%%%%%%%%%

For a period-\(k\) orbit

\[
x_1,\ldots,x_k,
\]

the orbit is locally attracting if

\[
\boxed{
\left|
\prod_{j=1}^{k}
f'(x_j)
\right|
<
1.
}
\]

This is equivalent to requiring

\[
\boxed{
|
(f^k)'(x_1)
|<1.
}
\]

%%%%%%%%%%%%%%%%%%%%%%%%%%%%%%%%%%%%%%%%%%%%%%%%%%%%%%%%%%%%
\subsection{Logistic Map}
\label{subsec:logistic-map}
%%%%%%%%%%%%%%%%%%%%%%%%%%%%%%%%%%%%%%%%%%%%%%%%%%%%%%%%%%%%

The logistic map is

\[
\boxed{
x_{n+1}
=
rx_n(1-x_n),
}
\]

often studied for

\[
\boxed{
0\leq x_n\leq1.
}
\]

Despite its simple formula, it exhibits fixed points, periodic cycles, and
chaos.

%%%%%%%%%%%%%%%%%%%%%%%%%%%%%%%%%%%%%%%%%%%%%%%%%%%%%%%%%%%%
\subsection{Logistic Map Fixed Points}
\label{subsec:logistic-map-fixed-points}
%%%%%%%%%%%%%%%%%%%%%%%%%%%%%%%%%%%%%%%%%%%%%%%%%%%%%%%%%%%%

Fixed points satisfy

\[
x
=
rx(1-x).
\]

Therefore,

\[
x
[
r(1-x)-1
]
=
0.
\]

Thus

\[
\boxed{
x^*_1=0,
}
\]

and for \(r\neq0\),

\[
\boxed{
x^*_2
=
1-\frac1r.
}
\]

%%%%%%%%%%%%%%%%%%%%%%%%%%%%%%%%%%%%%%%%%%%%%%%%%%%%%%%%%%%%
\subsection{Stability of \(x=0\)}
\label{subsec:logistic-map-zero-stability}
%%%%%%%%%%%%%%%%%%%%%%%%%%%%%%%%%%%%%%%%%%%%%%%%%%%%%%%%%%%%

For

\[
f(x)=rx(1-x),
\]

\[
\boxed{
f'(x)
=
r(1-2x).
}
\]

At \(x=0\),

\[
f'(0)=r.
\]

Thus \(x=0\) is attracting when

\[
\boxed{
|r|<1.
}
\]

In the common parameter range \(r\geq0\), this means

\[
0\leq r<1.
\]

%%%%%%%%%%%%%%%%%%%%%%%%%%%%%%%%%%%%%%%%%%%%%%%%%%%%%%%%%%%%
\subsection{Stability of the Nonzero Fixed Point}
\label{subsec:logistic-nonzero-stability}
%%%%%%%%%%%%%%%%%%%%%%%%%%%%%%%%%%%%%%%%%%%%%%%%%%%%%%%%%%%%

At

\[
x^*
=
1-\frac1r,
\]

\[
f'(x^*)
=
r
\left[
1
-
2
\left(
1-\frac1r
\right)
\right].
\]

Simplify:

\[
\boxed{
f'(x^*)
=
2-r.
}
\]

The fixed point is attracting when

\[
|2-r|<1.
\]

Therefore,

\[
\boxed{
1<r<3.
}
\]

%%%%%%%%%%%%%%%%%%%%%%%%%%%%%%%%%%%%%%%%%%%%%%%%%%%%%%%%%%%%
\subsection{Period Doubling}
\label{subsec:period-doubling}
%%%%%%%%%%%%%%%%%%%%%%%%%%%%%%%%%%%%%%%%%%%%%%%%%%%%%%%%%%%%

As the logistic-map parameter passes through

\[
r=3,
\]

the nonzero fixed point loses stability and a stable period-two orbit
appears.

Further parameter changes produce:

\[
\boxed{
2
\rightarrow
4
\rightarrow
8
\rightarrow
16
\rightarrow
\cdots
}
\]

periodic cycles.

This is called a period-doubling cascade.

%%%%%%%%%%%%%%%%%%%%%%%%%%%%%%%%%%%%%%%%%%%%%%%%%%%%%%%%%%%%
\subsection{Chaos}
\label{subsec:dynamical-chaos}
%%%%%%%%%%%%%%%%%%%%%%%%%%%%%%%%%%%%%%%%%%%%%%%%%%%%%%%%%%%%

Chaos is deterministic behavior that is:

\[
\boxed{
\text{aperiodic},
}
\]

\[
\boxed{
\text{bounded in many systems},
}
\]

and

\[
\boxed{
\text{highly sensitive to initial conditions}.
}
\]

Chaotic behavior is not the same as randomness.

The evolution rule remains deterministic.

%%%%%%%%%%%%%%%%%%%%%%%%%%%%%%%%%%%%%%%%%%%%%%%%%%%%%%%%%%%%
\subsection{Sensitive Dependence on Initial Conditions}
\label{subsec:sensitive-dependence}
%%%%%%%%%%%%%%%%%%%%%%%%%%%%%%%%%%%%%%%%%%%%%%%%%%%%%%%%%%%%

Two nearby initial states may separate rapidly.

If

\[
\delta_0
\]

is a small initial separation, one may observe approximately

\[
\boxed{
\delta_n
\approx
\delta_0e^{\lambda n}.
}
\]

When

\[
\lambda>0,
\]

small errors can grow exponentially.

%%%%%%%%%%%%%%%%%%%%%%%%%%%%%%%%%%%%%%%%%%%%%%%%%%%%%%%%%%%%
\subsection{Lyapunov Exponent for a Map}
\label{subsec:lyapunov-exponent-map}
%%%%%%%%%%%%%%%%%%%%%%%%%%%%%%%%%%%%%%%%%%%%%%%%%%%%%%%%%%%%

For

\[
x_{n+1}=f(x_n),
\]

a Lyapunov exponent may be defined as

\[
\boxed{
\lambda
=
\lim_{N\to\infty}
\frac1N
\sum_{n=0}^{N-1}
\log
|
f'(x_n)
|,
}
\]

when the limit exists.

If

\[
\boxed{
\lambda>0,
}
\]

nearby trajectories separate exponentially on average.

This is a major signature of chaos.

%%%%%%%%%%%%%%%%%%%%%%%%%%%%%%%%%%%%%%%%%%%%%%%%%%%%%%%%%%%%
\subsection{Predictability Horizon}
\label{subsec:predictability-horizon}
%%%%%%%%%%%%%%%%%%%%%%%%%%%%%%%%%%%%%%%%%%%%%%%%%%%%%%%%%%%%

Suppose initial uncertainty is

\[
\delta_0
\]

and acceptable prediction error is

\[
\Delta.
\]

If error grows as

\[
\delta(t)
\approx
\delta_0e^{\lambda t},
\]

then the approximate predictability horizon satisfies

\[
\boxed{
t
\approx
\frac1\lambda
\log
\left(
\frac{\Delta}{\delta_0}
\right).
}
\]

Thus chaotic systems can be deterministic yet have limited long-term
predictability.

%%%%%%%%%%%%%%%%%%%%%%%%%%%%%%%%%%%%%%%%%%%%%%%%%%%%%%%%%%%%
\subsection{Attractors}
\label{subsec:attractors}
%%%%%%%%%%%%%%%%%%%%%%%%%%%%%%%%%%%%%%%%%%%%%%%%%%%%%%%%%%%%

An attractor is an invariant set toward which nearby trajectories evolve.

Examples include:

\[
\boxed{
\text{stable equilibrium},
}
\]

\[
\boxed{
\text{stable limit cycle},
}
\]

and

\[
\boxed{
\text{chaotic attractor}.
}
\]

The basin of attraction consists of initial states converging toward the
attractor.

%%%%%%%%%%%%%%%%%%%%%%%%%%%%%%%%%%%%%%%%%%%%%%%%%%%%%%%%%%%%
\subsection{Strange Attractors Preview}
\label{subsec:strange-attractors-preview}
%%%%%%%%%%%%%%%%%%%%%%%%%%%%%%%%%%%%%%%%%%%%%%%%%%%%%%%%%%%%

A strange attractor is a chaotic attractor with complicated geometric
structure.

It may exhibit:

\[
\boxed{
\text{sensitive dependence},
}
\]

\[
\boxed{
\text{aperiodic trajectories},
}
\]

and

\[
\boxed{
\text{fractal-like geometry}.
}
\]

Examples arise in nonlinear systems of dimension three or higher.

%%%%%%%%%%%%%%%%%%%%%%%%%%%%%%%%%%%%%%%%%%%%%%%%%%%%%%%%%%%%
\subsection{Lorenz System Preview}
\label{subsec:lorenz-system-preview}
%%%%%%%%%%%%%%%%%%%%%%%%%%%%%%%%%%%%%%%%%%%%%%%%%%%%%%%%%%%%

A classical three-dimensional system is

\[
\boxed{
\begin{cases}
x'
=
\sigma(y-x),
\\
y'
=
x(\rho-z)-y,
\\
z'
=
xy-\beta z.
\end{cases}
}
\]

For certain parameter values, the Lorenz system exhibits chaotic
behavior.

Its equations are deterministic, but trajectories can be extremely
sensitive to initial conditions.

%%%%%%%%%%%%%%%%%%%%%%%%%%%%%%%%%%%%%%%%%%%%%%%%%%%%%%%%%%%%
\subsection{Deterministic Chaos Versus Randomness}
\label{subsec:chaos-versus-randomness}
%%%%%%%%%%%%%%%%%%%%%%%%%%%%%%%%%%%%%%%%%%%%%%%%%%%%%%%%%%%%

Random systems involve probabilistic evolution.

Chaotic deterministic systems use fixed deterministic equations.

However, both can produce complicated irregular-looking time series.

Distinguishing them from finite noisy data can be difficult.

%%%%%%%%%%%%%%%%%%%%%%%%%%%%%%%%%%%%%%%%%%%%%%%%%%%%%%%%%%%%
\subsection{Structural Stability Preview}
\label{subsec:structural-stability-preview}
%%%%%%%%%%%%%%%%%%%%%%%%%%%%%%%%%%%%%%%%%%%%%%%%%%%%%%%%%%%%

A dynamical behavior is structurally stable if sufficiently small changes
to the governing system do not change its qualitative phase portrait.

This is different from stability of an equilibrium.

Structural stability concerns robustness of the entire dynamical
structure.

%%%%%%%%%%%%%%%%%%%%%%%%%%%%%%%%%%%%%%%%%%%%%%%%%%%%%%%%%%%%
\subsection{Separatrices}
\label{subsec:separatrices}
%%%%%%%%%%%%%%%%%%%%%%%%%%%%%%%%%%%%%%%%%%%%%%%%%%%%%%%%%%%%

A separatrix is a trajectory or invariant curve that separates regions
with qualitatively different behavior.

Stable and unstable manifolds of saddle points frequently act as
separatrices.

They can divide different basins of attraction.

%%%%%%%%%%%%%%%%%%%%%%%%%%%%%%%%%%%%%%%%%%%%%%%%%%%%%%%%%%%%
\subsection{Poincaré Sections Preview}
\label{subsec:poincare-sections-preview}
%%%%%%%%%%%%%%%%%%%%%%%%%%%%%%%%%%%%%%%%%%%%%%%%%%%%%%%%%%%%

For higher-dimensional continuous systems, one can study intersections of
trajectories with a lower-dimensional surface.

The resulting return map is called a Poincaré map.

It converts aspects of continuous-time dynamics into a discrete
dynamical system.

This is useful for periodic-orbit and chaos analysis.

%%%%%%%%%%%%%%%%%%%%%%%%%%%%%%%%%%%%%%%%%%%%%%%%%%%%%%%%%%%%
\subsection{Return Maps}
\label{subsec:return-maps}
%%%%%%%%%%%%%%%%%%%%%%%%%%%%%%%%%%%%%%%%%%%%%%%%%%%%%%%%%%%%

Suppose a trajectory intersects a chosen section repeatedly.

Let

\[
x_n
\]

describe the \(n\)-th intersection.

A return map has form

\[
\boxed{
x_{n+1}
=
P(x_n).
}
\]

Fixed points of the return map correspond to periodic orbits of the
continuous system.

%%%%%%%%%%%%%%%%%%%%%%%%%%%%%%%%%%%%%%%%%%%%%%%%%%%%%%%%%%%%
\subsection{Dissipative Systems}
\label{subsec:dissipative-systems}
%%%%%%%%%%%%%%%%%%%%%%%%%%%%%%%%%%%%%%%%%%%%%%%%%%%%%%%%%%%%

A system is dissipative when phase-space volumes contract over time.

For

\[
\mathbf{x}'
=
\mathbf{f}(\mathbf{x}),
\]

the divergence

\[
\boxed{
\nabla\cdot\mathbf{f}
}
\]

measures local volume expansion or contraction.

Negative divergence often indicates local phase-volume contraction.

%%%%%%%%%%%%%%%%%%%%%%%%%%%%%%%%%%%%%%%%%%%%%%%%%%%%%%%%%%%%
\subsection{Divergence and Area Change}
\label{subsec:divergence-area-change}
%%%%%%%%%%%%%%%%%%%%%%%%%%%%%%%%%%%%%%%%%%%%%%%%%%%%%%%%%%%%

For a planar system

\[
\begin{cases}
x'=f(x,y),\\
y'=g(x,y),
\end{cases}
\]

the divergence is

\[
\boxed{
\frac{\partial f}{\partial x}
+
\frac{\partial g}{\partial y}.
}
\]

If this is strictly negative throughout a region, nearby phase-space area
tends to contract locally.

%%%%%%%%%%%%%%%%%%%%%%%%%%%%%%%%%%%%%%%%%%%%%%%%%%%%%%%%%%%%
\subsection{Bendixson--Dulac Criterion Preview}
\label{subsec:bendixson-dulac-preview}
%%%%%%%%%%%%%%%%%%%%%%%%%%%%%%%%%%%%%%%%%%%%%%%%%%%%%%%%%%%%

If the divergence has one strict sign throughout a simply connected
region under suitable smoothness assumptions, the planar system cannot
have a periodic orbit contained entirely within that region.

A more general version uses a Dulac function.

This provides a way to rule out limit cycles without solving the system.

%%%%%%%%%%%%%%%%%%%%%%%%%%%%%%%%%%%%%%%%%%%%%%%%%%%%%%%%%%%%
\subsection{Gradient Systems}
\label{subsec:gradient-systems}
%%%%%%%%%%%%%%%%%%%%%%%%%%%%%%%%%%%%%%%%%%%%%%%%%%%%%%%%%%%%

A gradient system has the form

\[
\boxed{
\mathbf{x}'
=
-\nabla V(\mathbf{x}).
}
\]

Then

\[
\frac{dV}{dt}
=
\nabla V^T
\mathbf{x}'.
\]

Thus,

\[
\boxed{
\frac{dV}{dt}
=
-
\|
\nabla V
\|^2
\leq
0.
}
\]

Therefore trajectories move toward lower values of the potential \(V\).

%%%%%%%%%%%%%%%%%%%%%%%%%%%%%%%%%%%%%%%%%%%%%%%%%%%%%%%%%%%%
\subsection{Equilibria of Gradient Systems}
\label{subsec:gradient-system-equilibria}
%%%%%%%%%%%%%%%%%%%%%%%%%%%%%%%%%%%%%%%%%%%%%%%%%%%%%%%%%%%%

Equilibria satisfy

\[
\boxed{
\nabla V(\mathbf{x}^*)=0.
}
\]

Local minima of \(V\) are natural candidates for stable equilibria.

Gradient systems cannot have nonconstant periodic orbits because \(V\)
would have to decrease and return to its original value.

%%%%%%%%%%%%%%%%%%%%%%%%%%%%%%%%%%%%%%%%%%%%%%%%%%%%%%%%%%%%
\subsection{Hamiltonian Systems Preview}
\label{subsec:hamiltonian-systems-dynamics}
%%%%%%%%%%%%%%%%%%%%%%%%%%%%%%%%%%%%%%%%%%%%%%%%%%%%%%%%%%%%

A planar Hamiltonian system has form

\[
\boxed{
\dot q
=
\frac{\partial H}{\partial p},
}
\]

\[
\boxed{
\dot p
=
-
\frac{\partial H}{\partial q}.
}
\]

The Hamiltonian \(H(q,p)\) is conserved:

\[
\boxed{
\frac{dH}{dt}=0.
}
\]

Thus trajectories lie on level curves of \(H\).

%%%%%%%%%%%%%%%%%%%%%%%%%%%%%%%%%%%%%%%%%%%%%%%%%%%%%%%%%%%%
\subsection{Continuous Versus Discrete Stability}
\label{subsec:continuous-discrete-stability}
%%%%%%%%%%%%%%%%%%%%%%%%%%%%%%%%%%%%%%%%%%%%%%%%%%%%%%%%%%%%

For the continuous linear scalar system

\[
x'=\lambda x,
\]

the equilibrium is asymptotically stable when

\[
\boxed{
\operatorname{Re}(\lambda)<0.
}
\]

For the discrete scalar system

\[
x_{n+1}
=
\lambda x_n,
\]

the equilibrium is asymptotically stable when

\[
\boxed{
|\lambda|<1.
}
\]

Thus the stability region differs between continuous and discrete time.

%%%%%%%%%%%%%%%%%%%%%%%%%%%%%%%%%%%%%%%%%%%%%%%%%%%%%%%%%%%%
\subsection{Numerical Exploration of Dynamical Systems}
\label{subsec:numerical-dynamical-systems}
%%%%%%%%%%%%%%%%%%%%%%%%%%%%%%%%%%%%%%%%%%%%%%%%%%%%%%%%%%%%

Many nonlinear systems have no elementary closed-form solution.

Numerical integration allows one to approximate trajectories.

A useful workflow includes:

\[
\boxed{
\text{solve equilibria analytically},
}
\]

\[
\boxed{
\text{analyze the Jacobian},
}
\]

\[
\boxed{
\text{plot nullclines},
}
\]

\[
\boxed{
\text{integrate representative trajectories},
}
\]

and

\[
\boxed{
\text{compare multiple initial conditions}.
}
\]

%%%%%%%%%%%%%%%%%%%%%%%%%%%%%%%%%%%%%%%%%%%%%%%%%%%%%%%%%%%%
\subsection{Numerical Artifacts}
\label{subsec:dynamical-numerical-artifacts}
%%%%%%%%%%%%%%%%%%%%%%%%%%%%%%%%%%%%%%%%%%%%%%%%%%%%%%%%%%%%

Numerical approximations can create behavior that is not present in the
true system.

Examples include:

\[
\boxed{
\text{artificial instability},
}
\]

\[
\boxed{
\text{false damping},
}
\]

or

\[
\boxed{
\text{spurious oscillation}.
}
\]

Step size and solver choice therefore matter.

%%%%%%%%%%%%%%%%%%%%%%%%%%%%%%%%%%%%%%%%%%%%%%%%%%%%%%%%%%%%
\subsection{Dynamical Systems Workflow}
\label{subsec:dynamical-systems-workflow}
%%%%%%%%%%%%%%%%%%%%%%%%%%%%%%%%%%%%%%%%%%%%%%%%%%%%%%%%%%%%

A practical continuous-system workflow is:

\begin{enumerate}
    \item write the system in autonomous first-order form;
    \item determine the physically relevant phase space;
    \item find all equilibria;
    \item compute nullclines;
    \item determine flow directions between nullclines;
    \item calculate the Jacobian;
    \item evaluate the Jacobian at each equilibrium;
    \item compute eigenvalues and eigenvectors;
    \item classify hyperbolic equilibria;
    \item use nonlinear analysis when linearization is inconclusive;
    \item identify invariant sets;
    \item look for conserved quantities or Lyapunov functions;
    \item investigate periodic behavior;
    \item vary parameters and identify bifurcations;
    \item use numerical trajectories to support qualitative analysis.
\end{enumerate}

%%%%%%%%%%%%%%%%%%%%%%%%%%%%%%%%%%%%%%%%%%%%%%%%%%%%%%%%%%%%
\subsection{Discrete-System Workflow}
\label{subsec:discrete-system-workflow}
%%%%%%%%%%%%%%%%%%%%%%%%%%%%%%%%%%%%%%%%%%%%%%%%%%%%%%%%%%%%

For

\[
x_{n+1}=f(x_n),
\]

a useful workflow is:

\begin{enumerate}
    \item determine the relevant state interval;
    \item solve \(f(x)=x\) for fixed points;
    \item compute \(f'(x^*)\);
    \item classify fixed-point stability;
    \item construct a cobweb diagram;
    \item search for periodic points using \(f^k(x)=x\);
    \item vary parameters;
    \item construct a bifurcation diagram;
    \item compute or estimate Lyapunov exponents when chaos is suspected.
\end{enumerate}

%%%%%%%%%%%%%%%%%%%%%%%%%%%%%%%%%%%%%%%%%%%%%%%%%%%%%%%%%%%%
\subsection{Common Errors}
\label{subsec:dynamical-systems-common-errors}
%%%%%%%%%%%%%%%%%%%%%%%%%%%%%%%%%%%%%%%%%%%%%%%%%%%%%%%%%%%%

\subsubsection{Confusing Equilibrium with Zero}

An equilibrium need not occur at

\[
\mathbf{x}=\mathbf{0}.
\]

It is any point satisfying

\[
\mathbf{f}(\mathbf{x}^*)=\mathbf{0}.
\]

%%%%%%%%%%%%%%%%%%%%%%%%%%%%%%%%%%%%%%%%%%%%%%%%%%%%%%%%%%%%
\subsubsection{Classifying Stability Without Checking All Eigenvalues}

For multidimensional systems, one unstable eigenvalue can make the entire
equilibrium unstable.

%%%%%%%%%%%%%%%%%%%%%%%%%%%%%%%%%%%%%%%%%%%%%%%%%%%%%%%%%%%%
\subsubsection{Assuming Negative Trace Always Implies Stability}

For a planar system, negative trace alone is insufficient.

If

\[
\det(A)<0,
\]

the equilibrium is a saddle and therefore unstable.

%%%%%%%%%%%%%%%%%%%%%%%%%%%%%%%%%%%%%%%%%%%%%%%%%%%%%%%%%%%%
\subsubsection{Ignoring the Determinant}

The trace and determinant must be considered together.

The determinant detects saddle behavior through

\[
\boxed{
\Delta<0.
}
\]

%%%%%%%%%%%%%%%%%%%%%%%%%%%%%%%%%%%%%%%%%%%%%%%%%%%%%%%%%%%%
\subsubsection{Using Linearization at a Nonhyperbolic Equilibrium as Final Proof}

If an eigenvalue has zero real part, nonlinear terms can determine
stability.

Linearization may be inconclusive.

%%%%%%%%%%%%%%%%%%%%%%%%%%%%%%%%%%%%%%%%%%%%%%%%%%%%%%%%%%%%
\subsubsection{Confusing a Center with an Asymptotically Stable Equilibrium}

A linear center has nearby closed trajectories.

They remain nearby but do not approach the equilibrium.

Thus a center is stable but not asymptotically stable.

%%%%%%%%%%%%%%%%%%%%%%%%%%%%%%%%%%%%%%%%%%%%%%%%%%%%%%%%%%%%
\subsubsection{Calling Every Closed Orbit a Limit Cycle}

A limit cycle must be isolated.

In conservative systems, a continuous family of closed energy curves is
not a family of isolated limit cycles.

%%%%%%%%%%%%%%%%%%%%%%%%%%%%%%%%%%%%%%%%%%%%%%%%%%%%%%%%%%%%
\subsubsection{Assuming Conservation Implies Attraction}

Conservative systems typically preserve energy and therefore do not
contract toward an attractor.

%%%%%%%%%%%%%%%%%%%%%%%%%%%%%%%%%%%%%%%%%%%%%%%%%%%%%%%%%%%%
\subsubsection{Assuming a Negative Lyapunov Derivative Must Be Strict Everywhere}

A negative semidefinite derivative can still support asymptotic stability
through additional arguments such as LaSalle's invariance principle.

%%%%%%%%%%%%%%%%%%%%%%%%%%%%%%%%%%%%%%%%%%%%%%%%%%%%%%%%%%%%
\subsubsection{Confusing Phase Space with Time Series Plots}

A phase portrait plots state variables against one another.

It is different from plotting one state variable against time.

%%%%%%%%%%%%%%%%%%%%%%%%%%%%%%%%%%%%%%%%%%%%%%%%%%%%%%%%%%%%
\subsubsection{Interpreting a Bifurcation as Numerical Error}

A genuine bifurcation is a qualitative change in the underlying
dynamical system as a parameter crosses a critical value.

%%%%%%%%%%%%%%%%%%%%%%%%%%%%%%%%%%%%%%%%%%%%%%%%%%%%%%%%%%%%
\subsubsection{Ignoring Parameter Dependence}

Stability conclusions can change when parameters change.

A system must be classified for the parameter regime being studied.

%%%%%%%%%%%%%%%%%%%%%%%%%%%%%%%%%%%%%%%%%%%%%%%%%%%%%%%%%%%%
\subsubsection{Confusing Chaos with Randomness}

Chaotic systems are deterministic.

Irregular behavior comes from nonlinear dynamics and sensitive dependence,
not from random forcing.

%%%%%%%%%%%%%%%%%%%%%%%%%%%%%%%%%%%%%%%%%%%%%%%%%%%%%%%%%%%%
\subsubsection{Assuming Sensitive Dependence Alone Proves Chaos}

A rigorous chaos definition generally requires additional structural
properties.

Positive Lyapunov exponent is a major diagnostic but is not the only
possible criterion.

%%%%%%%%%%%%%%%%%%%%%%%%%%%%%%%%%%%%%%%%%%%%%%%%%%%%%%%%%%%%
\subsubsection{Applying Continuous-Time Stability Rules to Discrete Maps}

Continuous systems use the sign of real parts of eigenvalues.

Discrete systems use eigenvalue magnitude relative to the unit circle.

%%%%%%%%%%%%%%%%%%%%%%%%%%%%%%%%%%%%%%%%%%%%%%%%%%%%%%%%%%%%
\subsubsection{Ignoring Numerical Solver Behavior}

A phase portrait generated numerically can be misleading if the step size
is too large or the method is unstable.

%%%%%%%%%%%%%%%%%%%%%%%%%%%%%%%%%%%%%%%%%%%%%%%%%%%%%%%%%%%%
\subsection{Applications}
\label{subsec:dynamical-systems-applications}
%%%%%%%%%%%%%%%%%%%%%%%%%%%%%%%%%%%%%%%%%%%%%%%%%%%%%%%%%%%%

\begin{applicationbox}

\textbf{Population Biology.}
Dynamical systems describe growth, competition, predator--prey
interaction, disease transmission, and ecosystem stability.

\medskip

\textbf{Epidemiology.}
Compartmental models such as SIR and SEIR systems use equilibria,
stability, thresholds, and bifurcations to describe epidemic behavior.

\medskip

\textbf{Physics.}
Pendulums, oscillators, Hamiltonian systems, celestial mechanics, and
nonlinear mechanical systems are naturally dynamical.

\medskip

\textbf{Engineering.}
Stability analysis is central to control systems, electrical networks,
feedback systems, and mechanical design.

\medskip

\textbf{Chemistry.}
Nonlinear chemical-reaction systems can exhibit steady states,
oscillations, and bifurcations.

\medskip

\textbf{Economics.}
Dynamic equilibrium, adjustment models, economic cycles, and iterative
market models can be studied using continuous or discrete systems.

\medskip

\textbf{Environmental Science.}
Coupled ecological, hydrological, and contaminant models often require
phase-plane and stability analysis.

\medskip

\textbf{Applied Mathematics.}
Dynamical systems provide a unifying language for nonlinear models across
science and engineering.

\end{applicationbox}

%%%%%%%%%%%%%%%%%%%%%%%%%%%%%%%%%%%%%%%%%%%%%%%%%%%%%%%%%%%%
\subsection{Exercises}
\label{subsec:dynamical-systems-exercises}
%%%%%%%%%%%%%%%%%%%%%%%%%%%%%%%%%%%%%%%%%%%%%%%%%%%%%%%%%%%%

\begin{exercise}[1]
Define a dynamical system.
\end{exercise}

\begin{exercise}[1]
Define phase space.
\end{exercise}

\begin{exercise}[1]
Define an equilibrium point.
\end{exercise}

\begin{exercise}[1]
Define Lyapunov stability.
\end{exercise}

\begin{exercise}[1]
Define asymptotic stability.
\end{exercise}

\begin{exercise}[1]
Define a basin of attraction.
\end{exercise}

\begin{exercise}[1]
Define a nullcline.
\end{exercise}

\begin{exercise}[1]
Define a limit cycle.
\end{exercise}

\begin{exercise}[1]
Define a bifurcation.
\end{exercise}

\begin{exercise}[1]
Define a fixed point of a discrete map.
\end{exercise}

\begin{exercise}[1]
Define chaos conceptually.
\end{exercise}

\begin{exercise}[2]
Find all equilibria of

\[
x'=x(x-2).
\]
\end{exercise}

\begin{exercise}[2]
Classify the equilibria of

\[
x'=x(x-2)
\]

using a phase line.
\end{exercise}

\begin{exercise}[2]
Use the derivative test to classify the equilibria of

\[
x'=x(3-x).
\]
\end{exercise}

\begin{exercise}[2]
Classify the origin for a linear system with eigenvalues

\[
-2,
\quad
-5.
\]
\end{exercise}

\begin{exercise}[2]
Classify the origin for eigenvalues

\[
3,
\quad
-1.
\]
\end{exercise}

\begin{exercise}[2]
Classify the origin for eigenvalues

\[
-2\pm4i.
\]
\end{exercise}

\begin{exercise}[2]
Classify the origin for eigenvalues

\[
2\pm i.
\]
\end{exercise}

\begin{exercise}[2]
For

\[
A
=
\begin{pmatrix}
1&2\\
3&4
\end{pmatrix},
\]

compute the trace and determinant.
\end{exercise}

\begin{exercise}[2]
Find the fixed points of

\[
x_{n+1}
=
2x_n(1-x_n).
\]
\end{exercise}

\begin{exercise}[2]
Classify the fixed point \(x=0\) for the map

\[
x_{n+1}
=
0.5x_n.
\]
\end{exercise}

\begin{exercise}[2]
Classify the fixed point \(x=0\) for

\[
x_{n+1}
=
1.3x_n.
\]
\end{exercise}

\begin{exercise}[3]
Construct the phase line for

\[
x'
=
x(1-x)(x-2).
\]
\end{exercise}

\begin{exercise}[3]
Classify every equilibrium of the previous equation.
\end{exercise}

\begin{exercise}[3]
Prove the one-dimensional derivative stability criterion using
linearization.
\end{exercise}

\begin{exercise}[3]
Explain the difference between stable and asymptotically stable.
\end{exercise}

\begin{exercise}[3]
Determine the basin of attraction for a simple scalar nonlinear system.
\end{exercise}

\begin{exercise}[3]
Classify a planar system using its trace and determinant.
\end{exercise}

\begin{exercise}[3]
Explain why

\[
\det(A)<0
\]

implies saddle behavior.
\end{exercise}

\begin{exercise}[3]
Find the nullclines of

\[
\begin{cases}
x'=x(1-x-y),\\
y'=y(2-x-y).
\end{cases}
\]
\end{exercise}

\begin{exercise}[3]
Find all equilibria of the previous system.
\end{exercise}

\begin{exercise}[3]
Compute the Jacobian of the previous system.
\end{exercise}

\begin{exercise}[3]
Classify the hyperbolic equilibria of the previous system.
\end{exercise}

\begin{exercise}[3]
Explain why linearization may fail at a center-type equilibrium.
\end{exercise}

\begin{exercise}[3]
Show that

\[
V=x^2+y^2
\]

is a Lyapunov function for

\[
x'=-3x,
\qquad
y'=-y.
\]
\end{exercise}

\begin{exercise}[3]
For a gradient system

\[
\mathbf{x}'
=
-\nabla V,
\]

prove that

\[
\dot V
=
-
\|
\nabla V
\|^2.
\]
\end{exercise}

\begin{exercise}[3]
Derive the conserved energy for

\[
x''+V'(x)=0.
\]
\end{exercise}

\begin{exercise}[3]
Analyze the equilibria of the undamped pendulum.
\end{exercise}

\begin{exercise}[3]
Find the predator--prey nullclines.
\end{exercise}

\begin{exercise}[3]
Find the predator--prey coexistence equilibrium.
\end{exercise}

\begin{exercise}[3]
Analyze the saddle-node normal form

\[
x'=\mu-x^2.
\]
\end{exercise}

\begin{exercise}[3]
Analyze the transcritical normal form

\[
x'=\mu x-x^2.
\]
\end{exercise}

\begin{exercise}[3]
Analyze the supercritical pitchfork normal form

\[
x'=\mu x-x^3.
\]
\end{exercise}

\begin{exercise}[3]
Find the logistic-map fixed points.
\end{exercise}

\begin{exercise}[3]
Derive the stability condition for the nonzero logistic-map fixed point.
\end{exercise}

\begin{exercise}[3]
Find all period-two points of a simple quadratic map by solving

\[
f^2(x)=x
\]

and excluding fixed points.
\end{exercise}

\begin{exercise}[3]
Explain how a cobweb diagram reveals attraction and repulsion.
\end{exercise}

\begin{exercise}[3]
Explain why a positive Lyapunov exponent limits long-range prediction.
\end{exercise}

\begin{exercise}[3]
Compare deterministic chaos and stochastic randomness.
\end{exercise}

\begin{exercise}[4]
Prove local asymptotic stability of a hyperbolic sink using
linearization theory.
\end{exercise}

\begin{exercise}[4]
Study the Hartman--Grobman theorem.
\end{exercise}

\begin{exercise}[4]
Study the stable-manifold theorem.
\end{exercise}

\begin{exercise}[4]
Study invariant manifolds near saddle points.
\end{exercise}

\begin{exercise}[4]
Apply LaSalle's invariance principle to a nonlinear system.
\end{exercise}

\begin{exercise}[4]
Study construction of quadratic Lyapunov functions for linear systems.
\end{exercise}

\begin{exercise}[4]
Derive the continuous-time Lyapunov matrix equation.
\end{exercise}

\begin{exercise}[4]
Study the Poincaré--Bendixson theorem.
\end{exercise}

\begin{exercise}[4]
Apply the Bendixson--Dulac criterion to rule out periodic orbits.
\end{exercise}

\begin{exercise}[4]
Study nonlinear pendulum separatrices.
\end{exercise}

\begin{exercise}[4]
Study conservative planar Hamiltonian systems.
\end{exercise}

\begin{exercise}[4]
Analyze coexistence and exclusion in a two-species competition model.
\end{exercise}

\begin{exercise}[4]
Study local bifurcation conditions for saddle-node bifurcations.
\end{exercise}

\begin{exercise}[4]
Study transcritical and pitchfork bifurcations using normal forms.
\end{exercise}

\begin{exercise}[4]
Study Hopf bifurcation theory.
\end{exercise}

\begin{exercise}[4]
Derive a normal form near a Hopf bifurcation.
\end{exercise}

\begin{exercise}[4]
Construct a logistic-map bifurcation diagram numerically.
\end{exercise}

\begin{exercise}[4]
Study the period-doubling cascade.
\end{exercise}

\begin{exercise}[4]
Investigate the Feigenbaum scaling constant.
\end{exercise}

\begin{exercise}[4]
Compute numerical Lyapunov exponents for the logistic map.
\end{exercise}

\begin{exercise}[4]
Study symbolic dynamics for chaotic maps.
\end{exercise}

\begin{exercise}[4]
Investigate Poincaré maps for periodically forced systems.
\end{exercise}

\begin{exercise}[4]
Study the Lorenz system and its strange attractor.
\end{exercise}

\begin{exercise}[4]
Investigate structural stability and bifurcation under parameter
perturbations.
\end{exercise}

%%%%%%%%%%%%%%%%%%%%%%%%%%%%%%%%%%%%%%%%%%%%%%%%%%%%%%%%%%%%
\subsection{Conceptual Summary}
\label{subsec:dynamical-systems-summary}
%%%%%%%%%%%%%%%%%%%%%%%%%%%%%%%%%%%%%%%%%%%%%%%%%%%%%%%%%%%%

A continuous autonomous dynamical system has form

\[
\boxed{
\mathbf{x}'
=
\mathbf{f}(\mathbf{x}).
}
\]

Equilibria satisfy

\[
\boxed{
\mathbf{f}(\mathbf{x}^*)=\mathbf{0}.
}
\]

For a scalar system

\[
x'=f(x),
\]

local stability can often be determined from

\[
\boxed{
f'(x^*).
}
\]

For a nonlinear multidimensional system, the Jacobian

\[
\boxed{
J(\mathbf{x}^*)
}
\]

provides the linearization.

Hyperbolic equilibria are classified by eigenvalues of the Jacobian.

Lyapunov functions provide a solution-free stability method through

\[
\boxed{
\dot V
=
\nabla V^T\mathbf{f}.
}
\]

A bifurcation is a qualitative change in dynamics as a parameter varies.

For discrete systems,

\[
\boxed{
x_{n+1}=f(x_n),
}
\]

fixed points satisfy

\[
\boxed{
f(x^*)=x^*.
}
\]

Their local stability is determined by

\[
\boxed{
|f'(x^*)|<1.
}
\]

Chaotic systems may exhibit

\[
\boxed{
\text{sensitive dependence on initial conditions},
}
\]

often associated with a positive Lyapunov exponent.

Dynamical systems therefore combine

\[
\boxed{
\text{differential equations}
+
\text{linear algebra}
+
\text{geometry}
+
\text{stability}
+
\text{nonlinear analysis}.
}
\]

%%%%%%%%%%%%%%%%%%%%%%%%%%%%%%%%%%%%%%%%%%%%%%%%%%%%%%%%%%%%
\subsection{Section Connections}
\label{subsec:dynamical-systems-connections}
%%%%%%%%%%%%%%%%%%%%%%%%%%%%%%%%%%%%%%%%%%%%%%%%%%%%%%%%%%%%

\chapterconnection{
Dynamical Systems shifts the emphasis from finding explicit solutions to
understanding the qualitative behavior of evolving systems. The ordinary
differential-equation methods of
Section~\ref{sec:ordinary-differential-equations} provide trajectories,
while eigenvalues, nullclines, Lyapunov functions, invariant sets, and
bifurcations explain long-term behavior. Discrete maps provide the
parallel framework for difference equations and reveal period doubling
and chaos. The next section, Partial Differential Equations, extends
differential-equation modeling to functions of several independent
variables and develops equations governing diffusion, waves, potentials,
and distributed physical systems.
}

%%%%%%%%%%%%%%%%%%%%%%%%%%%%%%%%%%%%%%%%%%%%%%%%%%%%%%%%%%%%
% SECTION SOURCE TRACKING
%%%%%%%%%%%%%%%%%%%%%%%%%%%%%%%%%%%%%%%%%%%%%%%%%%%%%%%%%%%%

%------------------------------------------------------------
% 9.3 Dynamical Systems
%------------------------------------------------------------
%
% Source basis:
%   - Standard qualitative dynamical-systems theory.
%   - Standard nonlinear ODE stability theory.
%   - Standard discrete dynamical systems and introductory chaos theory.
%
% Used for:
%   - state and phase space;
%   - trajectories, orbits, and flows;
%   - equilibria;
%   - scalar autonomous systems;
%   - phase-line analysis;
%   - Lyapunov, asymptotic, and exponential stability;
%   - basins of attraction;
%   - linear planar classification;
%   - trace-determinant analysis;
%   - nonlinear planar systems;
%   - nullclines;
%   - Jacobian linearization;
%   - hyperbolic equilibria;
%   - Hartman--Grobman preview;
%   - invariant sets;
%   - stable and unstable manifolds;
%   - Lyapunov functions;
%   - LaSalle preview;
%   - periodic orbits and limit cycles;
%   - conservative systems;
%   - pendulum dynamics;
%   - Poincare--Bendixson preview;
%   - predator--prey systems;
%   - competition models;
%   - saddle-node bifurcations;
%   - transcritical bifurcations;
%   - pitchfork bifurcations;
%   - Hopf bifurcations;
%   - bifurcation diagrams;
%   - discrete dynamical systems;
%   - fixed points and their stability;
%   - cobweb diagrams;
%   - periodic orbits;
%   - logistic map;
%   - period doubling;
%   - chaos;
%   - Lyapunov exponents;
%   - attractors and strange attractors;
%   - Lorenz-system preview;
%   - Poincare sections;
%   - dissipative and gradient systems;
%   - Hamiltonian-system preview;
%   - numerical dynamical analysis;
%   - applications and exercises.
%
% Notes:
%   - Partial Differential Equations is developed next in Section 9.4.
%   - Difference Equations later develop discrete systems more deeply.
%   - Numerical ODE and PDE methods are developed more deeply in
%     Chapter 11.
%
%%%%%%%%%%%%%%%%%%%%%%%%%%%%%%%%%%%%%%%%%%%%%%%%%%%%%%%%%%%%